% Bagian: computability
% Bab: recursive-functions
% Subbagian: other-recursions

\documentclass[../../../include/open-logic-section]{subfiles}

\begin{document}

\olfileid[id]{cmp}{rec}{ore}
\olsection{Bentuk Rekursi Lain}

Dengan menggunakan pemasangan dan pengodean barisan, kita dapat membenarkan
bentuk-bentuk rekursi primitif yang lebih eksotis (dan berguna). Sebagai
contoh, sering kali berguna untuk mendefinisikan dua fungsi secara serentak,
seperti dalam definisi berikut:
\begin{align*}
h_0(\vec x, 0) & = f_0(\vec x) \\
h_1(\vec x, 0) & = f_1(\vec x) \\
h_0(\vec x, y+1) & = g_0(\vec x, y, h_0(\vec x, y), h_1(\vec x, y)) \\
h_1(\vec x, y+1) & = g_1(\vec x, y, h_0(\vec x, y), h_1(\vec x, y)).
\end{align*}
Ini merupakan contoh \emph{rekursi serentak}. Cara lain yang berguna untuk
mendefinisikan fungsi ialah memberikan nilai $h(\vec x,y+1)$ berdasarkan
\emph{semua} nilai $h(\vec x,0)$, \dots,~$h(\vec x,y)$, seperti dalam
definisi berikut:
\begin{align*}
h(\vec x, 0) & = f(\vec x) \\
h(\vec x, y+1) & = g(\vec x, y, \tuple{h(\vec x, 0), \dots, h(\vec x, y)}).
\end{align*}
Skema berikut menangkap gagasan ini dengan lebih ringkas:
\[
h(\vec x, y) = g(\vec x, y, \tuple{h(\vec x, 0), \dots, h(\vec x, y-1)})
\]
dengan pengertian bahwa argumen terakhir bagi $g$ hanyalah barisan kosong
ketika $y=0$. Dalam kedua rumusan, gagasannya ialah bahwa saat menghitung
``langkah penerus'', fungsi $h$ dapat menggunakan seluruh barisan nilai yang
telah dihitung. Ini dikenal sebagai \emph{rekursi berdasarkan seluruh nilai
sebelumnya}. Sebagai contoh khusus, rekursi ini dapat digunakan untuk
membenarkan definisi berbentuk berikut:
\begin{align*}
h(\vec x, y) & = \begin{cases}
  g(\vec x, y, h(\vec x, k(\vec x, y))) & \text{jika $k(\vec x, y) < y$} \\
  f(\vec x) & \text{selain itu.}
\end{cases}
\end{align*}
Dengan kata lain, nilai $h$ pada $y$ dapat dihitung berdasarkan nilai $h$ pada
\emph{sebarang} nilai sebelumnya yang diberikan oleh~$k$.

\begin{prob}
  Definisikan fungsi sisa $r(x,y)$ melalui rekursi berdasarkan seluruh nilai
  sebelumnya. (Jika $x$, $y$ bilangan asli dan $y>0$, $r(x,y)$ merupakan
  bilangan yang lebih kecil daripada~$y$ sedemikian sehingga
  $x=z\times y+r(x,y)$ bagi suatu~$z$. Agar pasti, kita tetapkan bahwa jika
  $y=0$, maka $r(x,0)=0$.)
\end{prob}

Pikirkan bagaimana memperoleh fungsi-fungsi ini menggunakan rekursi primitif
biasa. Satu bentuk terakhir rekursi primitif lebih lentur karena kita
diperbolehkan mengubah \emph{parameter} (nilai tambahan) selama proses:
\begin{align*}
h(\vec x, 0) & = f(\vec x) \\
h(\vec x, y+1) & = g(\vec x, y, h(k(\vec x), y)).
\end{align*}
Bentuk ini juga dapat disimulasikan dengan rekursi primitif biasa. (Melakukan
hal itu tidak mudah. Sebagai petunjuk, coba uraikan komputasinya dengan tangan.)

\end{document}
